% Bao cao Pipes Puzzle - Greedy Best-First Search
% Bien dich: xelatex bao_cao.tex (chay hai lan de cap nhat muc luc)
\documentclass[12pt,a4paper]{article}

\usepackage{fontspec}
\usepackage{polyglossia}
\setdefaultlanguage{vietnamese}
\setmainfont{Times New Roman}
\usepackage{geometry}
\geometry{left=2.8cm,right=2.2cm,top=2.5cm,bottom=2.5cm}
\usepackage{setspace}
\onehalfspacing
\usepackage{amsmath,amssymb}
\usepackage{booktabs,longtable,tabularx,array,multirow}
\usepackage{graphicx}
\usepackage{xcolor}
\usepackage{float}
\usepackage{enumitem}
\usepackage{tikz}
\usetikzlibrary{arrows.meta,positioning,shapes.geometric,fit,calc}
\usepackage{listings}
\usepackage{fancyhdr}
\usepackage{hyperref}
\usepackage{microtype}
\usepackage{caption}
\usepackage{siunitx}
\sisetup{output-decimal-marker={,},group-separator={.},group-minimum-digits=4}

\definecolor{codebg}{RGB}{245,247,250}
\definecolor{codeblue}{RGB}{25,75,140}
\definecolor{codegreen}{RGB}{30,115,75}
\definecolor{darkblue}{RGB}{20,42,78}
\hypersetup{colorlinks=true,linkcolor=darkblue,urlcolor=codeblue,citecolor=darkblue}
\lstset{
  basicstyle=\ttfamily\small,
  backgroundcolor=\color{codebg},
  frame=single,
  rulecolor=\color{gray!35},
  breaklines=true,
  columns=fullflexible,
  keepspaces=true,
  showstringspaces=false,
  keywordstyle=\color{codeblue}\bfseries,
  commentstyle=\color{codegreen},
  numbers=left,
  numberstyle=\tiny\color{gray},
  xleftmargin=1.2em,
  framexleftmargin=1em
}
\pagestyle{fancy}
\fancyhf{}
\lhead{Pipes Puzzle}
\rhead{Greedy Best-First Search}
\cfoot{\thepage}

\newcommand{\code}[1]{\texttt{#1}}
\newcommand{\Hn}{\ensuremath{h(n)}}
\newcommand{\argmin}{\operatorname*{arg\,min}}

\begin{document}

\begin{titlepage}
\centering
\vspace*{1.1cm}
{\Large \textbf{TRƯỜNG ĐẠI HỌC BÁCH KHOA}}\\[0.25cm]
{\large \textbf{KHOA KHOA HỌC VÀ KỸ THUẬT MÁY TÍNH}}\\[1.4cm]
{\rule{\textwidth}{1.2pt}}\\[0.35cm]
{\huge \textbf{BÁO CÁO BÀI TẬP LỚN}}\\[0.35cm]
{\rule{\textwidth}{1.2pt}}\\[1.0cm]
{\LARGE \textbf{GIẢI PUZZLE PIPES BẰNG}}\\[0.2cm]
{\LARGE \textbf{GREEDY BEST-FIRST SEARCH}}\\[0.9cm]
\begin{tikzpicture}[scale=0.8]
  \fill[darkblue!12,rounded corners=5pt] (-3,-1.3) rectangle (3,1.3);
  \draw[darkblue,line width=4pt] (0,0)--(-1.8,0)--(-1.8,0.85);
  \draw[darkblue,line width=4pt] (0,0)--(1.8,0)--(1.8,-0.85);
  \draw[codegreen,line width=2pt] (-1.8,0.85)--(-1.8,0);
  \draw[codegreen,line width=2pt] (1.8,-0.85)--(1.8,0);
  \fill[darkblue] (0,0) circle (0.12);
  \fill[codegreen] (-1.8,0.85) circle (0.10);
  \fill[codegreen] (1.8,-0.85) circle (0.10);
\end{tikzpicture}\\[1.0cm]
\begin{tabular}{rl}
Sinh viên: & \rule{8cm}{0.4pt}\\[0.5cm]
MSSV: & \rule{8cm}{0.4pt}\\[0.5cm]
Lớp / Nhóm: & \rule{8cm}{0.4pt}\\[0.5cm]
Giảng viên: & \rule{8cm}{0.4pt}
\end{tabular}
\vfill
{\large TP. Hồ Chí Minh, năm 2026}
\end{titlepage}

\pagenumbering{roman}
\section*{Tóm tắt}
\addcontentsline{toc}{section}{Tóm tắt}
Báo cáo trình bày quá trình phân tích, mô hình hóa và hiện thực bài toán Pipes Puzzle bằng Python. Mục tiêu của trò chơi là xoay các đoạn ống trên một lưới để mọi đầu nối khớp nhau, không có đầu ống nào chĩa ra ngoài biên và toàn bộ mạng tạo thành một hệ thống liên thông. Hệ thống sử dụng Greedy Best-First Search, trong đó frontier được quản lý bởi hàng đợi ưu tiên và trạng thái có giá trị heuristic nhỏ nhất được mở rộng trước.

Ngoài thuật toán tìm kiếm, báo cáo mô tả bộ sinh puzzle procedural. Dense tạo một mạng liên thông ngẫu nhiên trên toàn lưới, bổ sung một số cạnh để hình thành nhánh và vòng lặp; Sparse giữ cây khung phủ toàn lưới và chọn các phép xoay làm đứt nhiều mối nối sẵn. Cả hai chế độ đều scramble có giới hạn nên bảo đảm bài có nghiệm. Giao diện Pygame được thiết kế theo dạng dashboard, cho phép chọn kích thước, xoay ống, chạy Greedy, theo dõi số trạng thái đã mở rộng theo thời gian thực và xem hoạt ảnh lời giải.

Kết quả đo trên bốn input mẫu cho thấy Greedy tìm được nghiệm trong tất cả trường hợp. Thời gian trung bình tăng từ khoảng \SI{0.0166}{s} ở lưới 3\(	imes\)3 lên \SI{0.2039}{s} ở lưới 6\(	imes\)6; bộ nhớ peak tương ứng tăng từ khoảng \SI{24.3}{KB} lên \SI{170.0}{KB}. Một phép đo BFS baseline ở giai đoạn khảo sát cho thấy frontier và bộ nhớ tăng nhanh, là lý do BFS không được giữ trong phiên bản cuối.

\textbf{Từ khóa:} Pipes Puzzle, Greedy Best-First Search, heuristic, Pygame, procedural generation, tracemalloc.

\tableofcontents
\listoftables
\clearpage
\pagenumbering{arabic}

\section{Giới thiệu}
\subsection{Bối cảnh}
Các bài toán tìm kiếm trên không gian trạng thái là một nội dung quan trọng của Trí tuệ nhân tạo. Một bài toán tìm kiếm thường gồm trạng thái ban đầu, tập phép biến đổi hợp lệ, kiểm tra đích và chiến lược chọn trạng thái tiếp theo. Pipes Puzzle có hình thức trực quan đơn giản nhưng thể hiện rõ các vấn đề của tìm kiếm: số trạng thái tăng theo số ô xoay được, nhiều trạng thái trung gian có cùng mức sai lệch, và nghiệm ngắn nhất không nhất thiết được tìm thấy bởi chiến lược tham lam.

Project ban đầu được tổ chức để so sánh tìm kiếm mù với tìm kiếm heuristic. Sau quá trình thử nghiệm, phạm vi bản cuối được thu gọn còn Greedy Best-First Search nhằm tập trung vào heuristic, sinh puzzle và trải nghiệm trực quan. BFS vẫn được dùng như baseline khảo sát trong báo cáo để giải thích quyết định thiết kế, nhưng không còn là chức năng chạy trong mã nguồn cuối.

\subsection{Mục tiêu}
Các mục tiêu chính của đề tài là:
\begin{enumerate}[label=\arabic*.]
  \item Mô hình hóa chính xác các loại ống, phép xoay và điều kiện thắng.
  \item Hiện thực Greedy Best-First Search có kiểm soát trạng thái đã thăm.
  \item Xây dựng heuristic phản ánh số đầu nối sai và mức rời rạc của mạng.
  \item Sinh nhiều puzzle ngẫu nhiên có cấu trúc dày, có nhánh/vòng và luôn có nghiệm.
  \item Thiết kế giao diện Pygame giúp người dùng quan sát board, thao tác thủ công và theo dõi tiến trình tìm kiếm.
  \item Đo thời gian, bộ nhớ, số node và frontier trên các input có kích thước khác nhau.
\end{enumerate}

\subsection{Phạm vi và giả thiết}
Mỗi ô chứa một trong các loại \code{END}, \code{STRAIGHT}, \code{ELBOW}, \code{TEE}, \code{CROSS} hoặc \code{EMPTY}. Một ô có thể xoay theo bốn góc 0, 90, 180 và 270 độ; các ống thẳng có hai hướng hình học khác nhau. Puzzle procedural có một ô nguồn cố định ở giữa; các ô ống khác đều có thể xoay, còn \code{CROSS} không sinh successor vì xoay không thay đổi kết nối.

Trong phạm vi báo cáo, chi phí của một phép xoay được xem là bằng nhau. Greedy ưu tiên tốc độ tìm nghiệm dựa trên \Hn và không cam kết nghiệm có ít phép xoay nhất.

\section{Phân tích bài toán}
\subsection{Luật của Pipes Puzzle}
Với mỗi cặp ô kề nhau, một đầu nối được xem là đúng khi cả hai phía cùng mở hoặc cùng đóng theo cùng hướng. Một đầu nối hướng ra ngoài biên luôn là lỗi. Ngoài kiểm tra cục bộ, trạng thái đích phải tạo ra một mạng liên thông: mọi ô có ống phải thuộc cùng thành phần liên thông.

Hình thức hóa điều kiện đích:
\begin{equation}
  \operatorname{goal}(s) \Longleftrightarrow
  \operatorname{mismatches}(s)=0
  \land \operatorname{connected\_to\_root}(s).
\end{equation}
Với puzzle được sinh bởi bộ sinh, các ô đều thuộc cùng một mạng bắt đầu từ nguồn trong trạng thái nghiệm. Hàm heuristic phạt theo số ô chưa nối được về nguồn để ưu tiên mở rộng một mạng duy nhất.

\subsection{Biểu diễn trạng thái}
Giả sử lưới có \(R\) hàng và \(C\) cột. Trạng thái được biểu diễn bằng tuple:
\begin{equation}
  s=(r_0,r_1,\ldots,r_{RC-1}),\qquad r_i\in\{0,1,2,3\}.
\end{equation}
Chỉ số một chiều của ô \((i,j)\) là \(iC+j\). Thông tin bất biến của ô, gồm loại ống và thuộc tính khóa, được lưu trong danh sách \code{Cell}; tuple trạng thái chỉ lưu góc xoay. Cách tách này giúp trạng thái có thể hash được để dùng trong \code{set} và \code{dict}.

Các hướng được mã hóa như sau:
\begin{center}
\begin{tabular}{c|cccc}
\toprule
Hướng & Bắc & Đông & Nam & Tây\\
\midrule
Mã & 0 & 1 & 2 & 3\\
\bottomrule
\end{tabular}
\end{center}

Mỗi loại ống có một mẫu kết nối chuẩn. Ví dụ, \code{STRAIGHT} ở rotation 0 nối Bắc--Nam, còn rotation 1 nối Đông--Tây; \code{ELBOW} ở rotation 0 nối Bắc--Đông. Hàm \code{rotated\_connections} quay mẫu chuẩn theo rotation.

\subsection{Successor và goal test}
Từ một trạng thái, successor được tạo bằng cách chọn một ô chưa khóa và chọn hướng xoay 90, 180 hoặc 270 độ đối với puzzle có nguồn cố định. Cách sinh trực tiếp này tránh buộc Greedy đi qua một trạng thái trung gian có heuristic xấu. Các trạng thái lặp được loại bằng tập \code{visited}. Với \code{CROSS} và \code{EMPTY}, việc xoay không tạo ra trạng thái có ý nghĩa nên không sinh successor.

Hàm \code{count\_mismatches} duyệt từng ô và bốn hướng. Mỗi cặp kết nối được đếm hai lần trong quá trình duyệt đối xứng, sau đó chia hai để thu được số lỗi theo cặp. Hàm \code{connectivity\_ok} dùng DFS trên đồ thị các kết nối đúng để kiểm tra các ô có ống có cùng thành phần hay không.

\section{Thuật toán Greedy Best-First Search}
\subsection{Ý tưởng}
Greedy Best-First Search chỉ dùng heuristic \Hn để chọn node tiếp theo:
\begin{equation}
  n^*=\argmin_{n\in OPEN} h(n).
\end{equation}
Frontier \code{OPEN} là một min-heap có phần tử dạng \((h(n),depth,tiebreaker,state)\). Độ sâu chỉ dùng phá hòa khi hai trạng thái có cùng heuristic; tiebreaker là bộ đếm tăng dần giúp thứ tự ổn định.

Thuật toán không cộng chi phí đã đi \(g(n)\), khác với A*. Vì vậy nó có thể đi theo một nhánh nhìn có vẻ tốt nhưng dài hơn nghiệm tối ưu. Đổi lại, thuật toán thường mở rộng ít node trên các puzzle Pipes có heuristic phản ánh trực tiếp số đầu nối lỗi.

\subsection{Heuristic}
Heuristic trong project là:
\begin{equation}
  h(s)=m(s)+2d(s),
\end{equation}
trong đó:
\begin{itemize}
  \item \(m(s)\) là số cặp đầu nối không khớp, bao gồm đầu nối chĩa ra ngoài biên;
  \item \(d(s)=1\) nếu đồ thị ống không liên thông và bằng 0 nếu liên thông.
\end{itemize}
Khoản phạt 2 cho thành phần rời rạc làm các trạng thái đã khớp cục bộ nhưng bị tách mạng kém hấp dẫn hơn. Heuristic bằng 0 là điều kiện đủ cho goal test trong mô hình hiện tại vì mọi lỗi biên và lỗi kề đều đã được tính vào \(m(s)\).

\subsection{Giả mã}
\begin{lstlisting}[language=Python,caption={Giả mã Greedy Best-First Search}]
GREEDY(puzzle, start):
    OPEN <- min_heap((h(start), 0, start))
    VISITED <- {start}
    PARENT[start] <- (None, "")
    while OPEN is not empty and generated <= MAX_NODES:
        state <- pop state with smallest h(state)
        expanded <- expanded + 1
        if goal(state):
            return reconstruct_path(PARENT, state)
        for (next_state, action) in successors(state):
            if next_state not in VISITED:
                VISITED.add(next_state)
                PARENT[next_state] <- (state, action)
                push (h(next_state), counter, next_state) into OPEN
    return failure
\end{lstlisting}

Độ phức tạp trong trường hợp xấu phụ thuộc vào branching factor \(b\) và độ sâu nghiệm \(d\). Không có bảo đảm tổng quát tốt hơn tìm kiếm heuristic tham lam; bộ nhớ chủ yếu đến từ \code{VISITED}, parent map và frontier. Tham số \code{max\_nodes} được dùng để ngăn chương trình chiếm tài nguyên không giới hạn.

\section{Sinh puzzle procedural}
\subsection{Mạng nghiệm ngẫu nhiên}
Bộ sinh mới không chọn từng loại ống hoàn toàn độc lập vì cách đó dễ tạo bài không có nghiệm. Quy trình gồm bốn bước:
\begin{enumerate}[label=\arabic*.]
  \item Chọn ô giữa làm nguồn và tạo cây khung bằng DFS ngẫu nhiên; sau bước này mọi ô trên lưới thuộc một mạng liên thông bắt đầu từ nguồn.
  \item Duyệt các cạnh lân cận chưa dùng và thêm một số cạnh với xác suất phụ thuộc kích thước. Các cạnh thêm tạo vòng lặp, nhánh và các nút bậc 3/4.
  \item Chuyển tập hướng kết nối của mỗi ô thành loại ống và rotation tương ứng. Bậc 1 tạo \code{END}, bậc 2 tạo \code{STRAIGHT} hoặc \code{ELBOW}, bậc 3 tạo \code{TEE}, bậc 4 tạo \code{CROSS}.
  \item Lưu rotation nghiệm và cố định ô nguồn ở giữa. Dense scramble một số ô bằng random walk giới hạn. Sparse giữ cây khung phủ toàn lưới, không thêm ô trống, rồi ưu tiên các phép xoay làm giảm nhiều nhất số cặp đầu nối đang khớp. Mỗi ô được xoay tối đa một lần trong quá trình scramble, với số ô thay đổi giới hạn khoảng một phần ba tổng số ô (tối thiểu bốn ô ở lưới nhỏ). Các ô khác vẫn có thể xoay để người chơi tự giải.
\end{enumerate}

Do scramble xuất phát từ trạng thái nghiệm, trạng thái ban đầu luôn có đường về nghiệm bằng cách đảo các phép xoay. Seed được hiển thị trong dashboard để việc tái lập một trường hợp có thể được bổ sung về sau.

\subsection{Mức độ khó}
Bốn tab tương ứng với các lưới 3\(\times\)3, 4\(\times\)4, 5\(\times\)5 và 6\(\times\)6. Cả Dense và Sparse đều giữ đủ ô ống trên toàn lưới. Dense giữ nhánh/vòng, còn Sparse chỉ dùng cây khung và chọn các phép xoay làm đứt nhiều mối nối sẵn để trạng thái ban đầu rời rạc hơn. Số ô scramble được giới hạn để kiểm soát độ khó. Khi kích thước tăng, nguồn trung tâm giữ vai trò neo kết nối cho toàn mạng.

\section{Kiến trúc và hiện thực}
\subsection{Cấu trúc module}
\begin{table}[H]
\centering
\caption{Các module chính của project}
\begin{tabularx}{\textwidth}{>{\ttfamily}l X}
\toprule
Module & Vai trò\\
\midrule
puzzle.py & Định nghĩa \code{Cell}, trạng thái, kết nối, successor, goal test và heuristic.\\
search.py & Hiện thực Greedy, parent map, heap, thống kê thời gian/bộ nhớ và callback tiến trình.\\
generator.py & Sinh mạng liên thông ngẫu nhiên từ nguồn trung tâm, thêm vòng/nhánh và scramble các ô.\\
main.py & Vòng lặp Pygame, board, dashboard, thao tác thủ công, worker search và hoạt ảnh lời giải.\\
benchmark.py & Chạy phép đo trên bốn fixture JSON và xuất các chỉ số.\\
\bottomrule
\end{tabularx}
\end{table}

\subsection{Luồng xử lý}
\begin{figure}[H]
\centering
\begin{tikzpicture}[
  node distance=10mm and 13mm,
  box/.style={draw=darkblue, rounded corners, fill=darkblue!7, align=center, minimum width=30mm, minimum height=11mm},
  arrow/.style={-{Latex[length=2mm]}, thick, darkblue}
]
\node[box] (input) {Người dùng\\chọn kích thước / N};
\node[box,right=of input] (gen) {Generator\\network + scramble};
\node[box,right=of gen] (model) {PipesPuzzle\\state + heuristic};
\node[box,below=of model] (search) {Greedy\\min-heap + visited};
\node[box,left=of search] (progress) {Progress callback\\expanded/generated};
\node[box,left=of progress] (ui) {Pygame dashboard\\và hoạt ảnh};
\draw[arrow] (input)--(gen);
\draw[arrow] (gen)--(model);
\draw[arrow] (model)--(search);
\draw[arrow] (search)--(progress);
\draw[arrow] (progress)--(ui);
\draw[arrow] (ui.north) |- (model.south);
\end{tikzpicture}
\caption{Luồng dữ liệu từ sinh puzzle đến hiển thị lời giải}
\end{figure}

\subsection{Giao diện Pygame}
Giao diện được tổ chức thành hai card ngang nhau. Card trái chứa lưới, tiêu đề, heuristic và legend trạng thái kết nối: xanh là đầu nối khớp, đỏ là lỗi kết nối; gốc nguồn được đánh dấu riêng. Card phải chứa trạng thái, bốn chỉ số search, nút chạy Greedy, step, pause, reset và cancel. Với puzzle nhỏ, tile block được căn giữa trong vùng board; legend được neo ở đáy card để layout không thay đổi khi đổi kích thước.

Các phím điều khiển chính là: `1--4` đổi kích thước, `N` tạo seed/layout mới, `G` chạy Greedy, `S` step, `Space` pause/continue, `R` reset, `C` cancel và click chuột để xoay ô chưa khóa.

\subsection{Hiển thị tiến trình theo thời gian thực}
Nếu tìm kiếm chạy đồng bộ, cửa sổ sẽ không vẽ được trong thời gian Greedy mở rộng frontier. Vì vậy phiên bản hiện tại dùng một worker nền. Main loop vẫn xử lý event và vẽ frame; worker gọi progress callback sau mỗi lần mở rộng. Callback cập nhật:
\begin{itemize}
  \item số node đã mở rộng;
  \item số node đã sinh;
  \item kích thước frontier hiện tại và frontier lớn nhất;
  \item thời gian từ lúc bắt đầu search.
\end{itemize}
Khi worker trả về \code{SearchResult}, main loop chuyển sang hoạt ảnh path. Cờ token và cancel check ngăn kết quả của worker cũ ghi đè puzzle mới khi người dùng nhấn reset hoặc đổi tab.

\section{Thiết kế thực nghiệm}
\subsection{Môi trường và quy trình}
Các phép đo được chạy trong môi trường:
\begin{itemize}
  \item Python 3.13.7, Pygame 2.6.1;
  \item Windows 10/11, chạy trong virtual environment \code{venv};
  \item \code{max\_nodes = 500000} cho Greedy fixture;
  \item \code{time.perf\_counter()} đo thời gian tường (wall-clock);
  \item \code{tracemalloc} đo peak memory do Python cấp phát trong search;
  \item mỗi fixture Greedy được chạy 5 lần, báo cáo lấy trung bình và độ lệch chuẩn.
\end{itemize}

Bốn fixture là các file JSON 3\(\times\)3 đến 6\(\times\)6 trong thư mục \code{puzzles}. Đây là input cố định dùng để so sánh lặp lại; các puzzle procedural tạo trong giao diện được kiểm tra riêng bằng nhiều seed.

\begin{table}[H]
\centering
\caption{Đặc trưng input dùng cho benchmark Greedy}
\begin{tabular}{lrrrr}
\toprule
Input & Kích thước & Số ô & Max node & Số lần chạy\\
\midrule
\code{tiny\_3x3} & 3\(\times\)3 & 9 & 500000 & 5\\
\code{easy\_4x4} & 4\(\times\)4 & 16 & 500000 & 5\\
\code{medium\_5x5} & 5\(\times\)5 & 25 & 500000 & 5\\
\code{hard\_6x6} & 6\(\times\)6 & 36 & 500000 & 5\\
\bottomrule
\end{tabular}
\end{table}

\section{Kết quả và phân tích}
\subsection{Greedy Best-First Search}
Bảng \ref{tab:greedy} sử dụng số liệu trung bình của năm lần chạy. \emph{Depth} là số phép xoay trong path do Greedy trả về; \emph{expanded} là số node đã pop khỏi heap; \emph{generated} là tổng số trạng thái đã sinh; \emph{frontier max} là kích thước lớn nhất của heap.

\begin{table}[H]
\centering
\caption{Thời gian và bộ nhớ của Greedy trên các input cố định}
\label{tab:greedy}
\small
\begin{tabular}{lrrrrrrr}
\toprule
Input & Found & Depth & Expanded & Generated & Frontier max & Time (s) & Peak (KB)\\
\midrule
Tiny 3x3 & Có & 7 & 15 & 100 & 86 & 0,0166 $\pm$ 0,0003 & 24,3 $\pm$ 9,8\\
Easy 4x4 & Có & 12 & 18 & 194 & 177 & 0,0514 $\pm$ 0,0011 & 39,0 $\pm$ 18,0\\
Medium 5x5 & Có & 17 & 22 & 331 & 310 & 0,1295 $\pm$ 0,0036 & 143,1 $\pm$ 7,0\\
Hard 6x6 & Có & 12 & 18 & 331 & 314 & 0,2039 $\pm$ 0,0058 & 170,0 $\pm$ 0,1\\
\bottomrule
\end{tabular}
\end{table}

Kết quả cho thấy thời gian tăng theo kích thước và số trạng thái sinh, nhưng không tăng đơn điệu theo depth vì Greedy không tối ưu path. Hard 6x6 có depth ngắn hơn Medium 5x5 trong fixture này nhưng frontier và peak memory vẫn cao hơn. Đây là lý do cần báo cáo đồng thời cả depth, node và memory thay vì chỉ nhìn thời gian.

Bộ nhớ peak của Tiny và Easy có độ lệch chuẩn tương đối lớn vì kích thước bài nhỏ, nên chi phí khởi tạo và cấp phát của Python chiếm tỷ trọng cao. Với Medium và Hard, phần frontier/visited lớn hơn nên số đo ổn định hơn. Cần hiểu peak của \code{tracemalloc} là bộ nhớ Python được theo dõi, không phải toàn bộ RSS của tiến trình Windows.

\subsection{BFS baseline trong giai đoạn khảo sát}
BFS đã được chạy ở giai đoạn đầu để làm đường cơ sở, trước khi loại khỏi phiên bản cuối. Để tránh một phép đo không kết thúc chiếm tài nguyên, các case Easy--Hard bị giới hạn ở 10000 trạng thái sinh. Tiny tìm được nghiệm trước giới hạn; các case lớn hơn chạm giới hạn nên cột Found là Không. Đây là bảng giải thích chi phí của BFS, không phải API còn lại trong project cuối.

\begin{table}[H]
\centering
\caption{BFS baseline với giới hạn 10000 trạng thái sinh}
\label{tab:bfs}
\small
\begin{tabular}{lrrrrrr}
\toprule
Input & Found & Expanded & Generated & Frontier max & Time (s) & Peak (KB)\\
\midrule
Tiny 3x3 & Có & 1810 & 3849 & 2040 & 0,6206 & 715,8\\
Easy 4x4 & Không & 3039 & 10001 & 6963 & 2,6353 & 2214,7\\
Medium 5x5 & Không & 2024 & 10001 & 7978 & 8,5529 & 2923,3\\
Hard 6x6 & Không & 1572 & 10001 & 8430 & 12,5006 & 3786,4\\
\bottomrule
\end{tabular}
\end{table}

BFS dùng queue FIFO nên không tận dụng mức độ sai lệch của trạng thái. Nó phải giữ nhiều trạng thái cùng độ sâu trong frontier, dẫn đến memory peak lớn hơn Greedy nhiều lần. Tiny là trường hợp BFS còn chấp nhận được; từ Easy trở lên, chỉ 10000 trạng thái đã cần nhiều bộ nhớ và chưa tìm được nghiệm trong phép đo giới hạn. Kết quả này hỗ trợ quyết định giữ Greedy trong bản chạy chính.

\subsection{Kiểm tra puzzle procedural}
Ngoài fixture cố định, bộ sinh được kiểm tra với nhiều seed 6\(\times\)6. Một nhóm năm seed đại diện cho kết quả: heuristic ban đầu nằm trong khoảng 27--35, Greedy đều tìm thấy nghiệm, depth nằm trong khoảng 36--46, expanded nằm trong khoảng 112--562. Điều này cho thấy layout procedural thay đổi thực sự nhưng vẫn nằm trong giới hạn phản hồi của giao diện.

\section{Đánh giá thiết kế}
\subsection{Ưu điểm}
\begin{enumerate}[label=\arabic*.]
  \item Heuristic đơn giản, dễ giải thích và gắn trực tiếp với luật nối ống.
  \item Parent map cho phép khôi phục path mà không lưu toàn bộ path trong mỗi node.
  \item Generator tạo mạng liên thông trước khi scramble nên không sinh bài vô nghiệm do cấu trúc.
  \item Seed và kích thước được hiển thị; nút `N` tạo layout mới, còn `R` giúp thử lại cùng layout.
  \item Worker và progress callback giúp số liệu trên dashboard tăng trong lúc tìm kiếm thay vì đóng băng giao diện.
  \item Tracemalloc, frontier, expanded và generated cung cấp góc nhìn đầy đủ hơn một phép đo thời gian đơn lẻ.
\end{enumerate}

\subsection{Hạn chế}
\begin{enumerate}[label=\arabic*.]
  \item Greedy không bảo đảm số phép xoay tối ưu và có thể bị dẫn sai bởi heuristic.
  \item Vì mọi ô (trừ nguồn và CROSS) đều có thể xoay, puzzle lớn vẫn có thể tạo frontier rộng; số bước scramble được giới hạn để giữ tốc độ phản hồi.
  \item Peak memory hiện là peak Python allocation từ tracemalloc; chưa đo GPU hoặc RSS toàn hệ thống.
  \item Worker thread giúp UI không treo nhưng cần token/cancel check để tránh race khi reset hoặc đổi puzzle.
  \item Bảng BFS baseline được đo bằng evaluator khảo sát với giới hạn node, vì BFS đã loại khỏi bản cuối.
\end{enumerate}

\subsection{Hướng phát triển}
Các hướng phát triển hợp lý gồm: bổ sung A* để so sánh trực tiếp trade-off giữa \(g(n)\) và \(h(n)\); cho phép người dùng nhập seed để tái lập bài; sinh difficulty dựa trên số phép xoay tối thiểu thay vì chỉ số scramble; ghi kết quả benchmark ra CSV tự động; và thêm biểu đồ thời gian/memory theo kích thước lưới.

\section{Kết luận}
Đề tài đã hiện thực đầy đủ một pipeline từ mô hình hóa Pipes Puzzle, tạo successor và goal test, xây dựng heuristic, tìm kiếm Greedy, sinh puzzle procedural đến giao diện Pygame. Greedy chọn trạng thái có \Hn thấp nhất bằng min-heap và sử dụng visited để tránh lặp. Bộ sinh mới tạo các mạng dày có nhánh/vòng, giúp mỗi lần chơi có layout khác nhau nhưng vẫn giữ nghiệm.

Các phép đo cho thấy Greedy phù hợp với mục tiêu tương tác: trên bốn fixture, thời gian và bộ nhớ thấp, nghiệm được tìm thấy nhanh; trong khi BFS baseline nhanh chóng tăng frontier và memory khi kích thước tăng. Vì Greedy không tối ưu đường đi, báo cáo đã phân tích đồng thời depth, số node, frontier, thời gian và peak memory. Giao diện cuối cùng cho phép quan sát các chỉ số này theo thời gian thực, làm rõ cách thuật toán hoạt động thay vì chỉ hiển thị trạng thái thắng/thua.

\appendix
\section{Hướng dẫn chạy và tái lập}
\begin{lstlisting}[language=bash,caption={Các lệnh chạy project}]
python -m venv venv
.\venv\Scripts\Activate.ps1
pip install -r requirements.txt
python main.py
python benchmark.py
\end{lstlisting}

Để biên dịch báo cáo bằng XeLaTeX:
\begin{lstlisting}[language=bash]
xelatex bao_cao.tex
xelatex bao_cao.tex
\end{lstlisting}

Các file quan trọng là \code{puzzle.py}, \code{search.py}, \code{generator.py}, \code{main.py} và \code{benchmark.py}. Khi thay đổi heuristic hoặc generator, cần chạy lại benchmark và ghi rõ seed, giới hạn node, số lần lặp và phiên bản Python trong phần thực nghiệm.

\section{Tài liệu tham khảo}
\begin{enumerate}
  \item Python Software Foundation, \emph{Python 3 Documentation: heapq, collections and threading}.
  \item Python Software Foundation, \emph{tracemalloc --- Trace memory allocations}.
  \item Pygame Community, \emph{Pygame Documentation}.
  \item Stuart Russell and Peter Norvig, \emph{Artificial Intelligence: A Modern Approach}, phần tìm kiếm heuristic.
\end{enumerate}

\end{document}
